\documentclass[tikz]{standalone}
\usepackage{amsmath}
\definecolor{bg}{HTML}{171D2B}
\definecolor{acc}{HTML}{F7B538}
\definecolor{mut}{HTML}{8A94A6}
\begin{document}
\pagecolor{bg}
\begin{tikzpicture}[line width=1.1pt, text=white]
% ---------- KIRI: kaleng dengan benang dililit ----------
% kaleng: silinder tampak depan
\fill[acc!25] (-1.6,-1.9) rectangle (1.6,1.9);
\fill[bg] (-1.6,-1.9) rectangle (1.6,1.9);
\draw[acc] (-1.6,1.7) arc (180:360:1.6 and 0.42);
\draw[acc] (-1.6,-1.7) ellipse (1.6 and 0.42);
\draw[acc] (-1.6,1.7) -- (-1.6,-1.7);
\draw[acc] (1.6,1.7) -- (1.6,-1.7);
% garis-garis lilitan benang (diagonal)
\foreach \i in {0,...,6} {
  \draw[white, line width=0.6pt] (-1.55,{1.7-\i*0.55}) -- (1.55,{1.15-\i*0.55});
  \draw[white, line width=0.6pt] (1.55,{1.7-\i*0.55}) -- (-1.55,{1.15-\i*0.55});
}
\draw[white, very thick, -{Stealth[length=3mm]}] (0,0) -- (0,-3.2) node[below, text=white] {$K$ (keliling) dari benang};
% ---------- KANAN: lingkaran antara dua buku ----------
\begin{scope}[shift={(6.6,0)}]
  % dua buku sejajar
  \fill[mut!30] (-2.6,-2.4) rectangle (2.6,-1.9);
  \fill[mut!30] (-2.6,1.9) rectangle (2.6,2.4);
  \draw[mut] (-2.6,-2.4) rectangle (2.6,-1.9);
  \draw[mut] (-2.6,1.9) rectangle (2.6,2.4);
  % lingkaran di antara buku
  \draw[acc, line width=1.6pt] (0,0) circle (1.9);
  \fill[white] (0,0) circle (1.6pt);
  % diameter
  \draw[white, -{Stealth[length=3mm]}, {Stealth[length=3mm]}-] (-1.9,0) -- (1.9,0) node[midway, above, text=white] {$d$ (diameter)};
  \node[below, text=white] at (0,-2.9) {$d$ = jarak antar buku};
\end{scope}
% hubungan
\node[text=white, font=\large] at (3.1,-3.9) {$\dfrac{K}{d} \approx 3{,}14$};
\end{tikzpicture}
\end{document}
